La criptografía está presente inadvertidamente en nuestro día a día: llamadas de teléfono, aplicaciones de mensajería como WhatsApp o Telegram, navegación segura por internet, la conexión Bluetooth entre dispositivos, pagar con tarjeta en un comercio, actualizaciones automáticas de teléfonos u ordenadores, la declaración de la renta, etc. La criptografía es la disciplina que se encarga de proporcionar seguridad a la hora de establecer estas comunicaciones. Más concretamente, es el estudio de técnicas matemáticas que garantizan la seguridad de la información en presencia de adversarios. Con ``seguridad'' se hace referencia a unos objetivos específicos derivados de la propia comunicación donde está implementada: al llamar por teléfono o enviar un mensaje por WhatsApp, la información se encripta de tal modo que si alguien lograra interceptarlos, no podría obtener información acerca de la comunicación; a la hora de navegar por internet, el navegador (cliente) necesita cerciorarse de que efectivamente ha establecido conexión con el sitio web deseado (servidor), y no con un tercero que esté suplantando la identidad del servidor, así como que la información que recibe es la enviada por el servidor, no habiendo sido manipulada en el proceso; a la hora de hacer la declaración de la renta, el documento queda ligado a quien lo firmó, de tal forma que se podría probar ante un tercero su autoría.

Por lo tanto, hay tantos objetivos de seguridad como necesidades surgen a la hora de proteger la información en un sistema de comunicación. Estos objetivos pueden agruparse en cuatro objetivos más generales:

\begin{itemize}
	\item \textbf{Confidencialidad}: Garantiza que la información solo sea accesible por aquellos autorizados.
	\item \textbf{Integridad}: Garantiza que los datos no hayan sido modificados durante la transmisión.
	\item \textbf{Autenticación}: Garantiza la identidad del emisor de los datos.
	\item \textbf{No repudio}: Garantiza que una entidad no pueda denegar acciones realizadas previamente.
\end{itemize}

Para alcanzar estos objetivos, la criptografía moderna se construye a partir de primitivas criptográficas: algoritmos básicos diseñados para resolver una única tarea de seguridad. Para cumplir varios objetivos a la vez o proporcionar una mayor funcionalidad, varias primitivas pueden combinarse junto con otros algoritmos para formar un sistema criptográfico o ``criptosistema''. Cuando dos o más partes necesitan intercambiar información para lograr uno o varios objetivos comunes, primitivas y criptosistemas se integran para formar un protocolo. Estos términos pueden diferir ligeramente de autor a autor, por lo que lo principal es \begin{Revision}entender\end{Revision} que las primitivas son los algoritmos básicos a partir de los cuales se construyen sistemas y protocolos que dan solución a uno o varios objetivos de seguridad. 

\begin{Revision}Este texto se centrará en la construcción de los cuerpos finitos, estructura algebraica en la cual se fundamentan varios pasos del algoritmo AES (una primitiva criptográfica cuyo propósito es proveer confidencialidad), así como mostrar cómo se comporta la aritmética en dicha estructura. Además, se explicará también el funcionamiento interno de AES, siendo una excelente oportunidad para aplicar los conceptos aprendidos sobre cuerpos finitos. Con este fin, el presente trabajo se divide en 3 capítulos. El primero contextualizará el estudio de la criptografía, explicando cómo se procesa la información y los tipos de sistemas criptográficos actuales, como los basados en clave simétrica y asimétrica. El segundo capítulo ocupa la mayor parte del texto, y está centrado en construir, de forma razonada y pedagógica, los cuerpos finitos: conjuntos con finitos elementos donde se pueden realizar las operaciones de suma, resta, multiplicación y división, además de entender cómo se comportan estas operaciones en dicha estructura. El tercer y último capítulo se centra en el algoritmo criptográfico AES, su estructura y funcionamiento. Se enfatiza, por tanto, que el objetivo de este trabajo es el estudio desde el punto de vista algebraico de AES. Quedan así fuera del ámbito de estudio los aspectos históricos de la criptografía y cómo se ha llegado a la criptografía moderna, el resto de criptosistemas y algoritmos de clave simétrica y pública, las funciones hash, y el estudio de vulnerabilidades de AES (criptoanálisis) y la filosofía detrás de su diseño.
\end{Revision}

